\documentclass[11pt,a4paper]{article}
\usepackage[margin=1in]{geometry}
\usepackage{amsmath,amssymb}
\usepackage[utf8]{inputenc}
\usepackage[T1]{fontenc}
\usepackage{graphicx}
\usepackage{booktabs}
\usepackage{multirow}
\usepackage{xcolor}
\usepackage{hyperref}
\usepackage[right]{lineno}
\usepackage{microtype}
\usepackage{natbib}
\usepackage{caption}
\usepackage{subcaption}
\usepackage{float}

\hypersetup{colorlinks=true, linkcolor=blue!60!black, citecolor=blue!60!black, urlcolor=blue!60!black}

\linenumbers

\graphicspath{{../results/celltype_landscape/quantitative_modulation_paper/figures/}}

\title{\textbf{A high-confidence causal bQTL atlas reveals the regulatory
architecture of natural gene expression variation in maize}}

\author{
Michael~Banf\textsuperscript{1} and Thomas~Hartwig\textsuperscript{2,3,4,*}
\\[6pt]
\small\textsuperscript{1}Perelyn, Munich, Germany\\
\small\textsuperscript{2}Heinrich Heine University D\"usseldorf, Institute for Molecular Physiology, D\"usseldorf, Germany\\
\small\textsuperscript{3}Max Planck Institute for Plant Breeding Research, Cologne, Germany\\
\small\textsuperscript{4}Cluster of Excellence on Plant Sciences (CEPLAS), D\"usseldorf, Germany\\[4pt]
\small *~Corresponding author: hartwit@hhu.de
}

\date{}

\begin{document}
\maketitle

% ============================================================
\begin{abstract}
% ============================================================

Binding quantitative trait loci (bQTL) identify sequence variants that alter
transcription factor (TF) occupancy, but linking individual variants to causal
gene expression changes remains challenging due to linkage disequilibrium and
motif redundancy. Here, we apply a six-step filtering pipeline to 147,942
maize bQTL from a 25-hybrid pan-cistrome to isolate 728 high-confidence
\emph{causal} variants---positions where a single nucleotide change at a
single bQTL per gene causally shifts allele-specific expression (ASE) across
multiple NAM hybrids (FDR\,$<$\,0.05). We demonstrate that within the actual
${\sim}$160\,bp MOA-seq binding peak, the disrupted TF family's motif is the
\emph{sole copy} in 47.8\% of bQTL (and 28.2\% of all family-at-variant instances),
with striking variation across families---BES1 (86\%), Nin-like (85\%), and
CAMTA (70\%) are almost always non-redundant, while C2H2, MYB, and ERF
(${\sim}$10\%) maintain multiple copies even within narrow peaks. Crucially,
32.5\% of variants simultaneously disrupt one TF family's motif while
creating another---with zero self-switches---revealing that causal bQTL
\emph{rewire} TF occupancy rather than simply ablate it. These 728 target
chaperones (8.7\%), and transcription factors (4.4\%)---rather than metabolic
enzymes, and are enriched for the pentose phosphate pathway (OR\,=\,3.37),
steroid biosynthesis (OR\,=\,3.60), and ubiquitin-mediated proteolysis
(OR\,=\,2.14). Clustering by disrupted TF motif reveals that ERF and MYB
each exclusively regulate 20 target genes, while all TF families converge on
the same functional targets (signal transduction, protein quality control),
consistent with chromatin accessibility---not TF identity---being the primary
gate for functional variation. We identify 34 transcription factor genes that
are themselves causal bQTL targets, creating regulatory cascades where a
single upstream motif disruption can propagate to 26--96 downstream genes.
This atlas provides a prioritized set of 728 single-variant regulatory
switches suitable for functional validation and precision breeding.

\end{abstract}

% ============================================================
\section{Introduction}
% ============================================================

A central question in quantitative genetics is how sequence variation at
cis-regulatory elements translates into phenotypic diversity. In maize, the
pan-cistrome study by Engelhorn et al.\ \citep{engelhorn2025} identified
147,942 binding quantitative trait loci (bQTL)---positions where at least one
of 25~F1 hybrids (B73 $\times$ diverse founders) shows allelic differences in
TF binding relative to the B73 reference. These bQTL collectively explained
the majority of heritable trait variation for over 70\% of 143 phenotypes
tested, establishing cis-regulatory variation as a dominant force in maize
biology.

However, not every bQTL is equally informative. Many bQTL are in linkage
disequilibrium (LD), cluster in the same promoter, or disrupt motifs that have
redundant copies nearby. To move from correlative association to mechanistic
understanding, we need to identify variants where we can confidently attribute
a causal chain: \emph{this specific nucleotide} disrupts \emph{this specific
TF binding site}, which shifts expression of \emph{this specific gene}.

Here, we describe a systematic filtering approach that distills 147,942 bQTL
down to 728 high-confidence causal variants. Each passes three criteria:
(1)~it is the \emph{only} significant bQTL within 2\,kb of its target gene's
TSS, eliminating LD ambiguity; (2)~it falls within an actual MOA-seq binding
peak, confirming the variant is in actively bound chromatin; and (3)~the
genotype at this position causally predicts allele-specific expression (ASE)
across 19 genotyped NAM hybrids (FDR\,$<$\,0.05). We then characterize these
728 genes and their disrupted TF motifs to reconstruct the regulatory
architecture of natural expression variation.


% ============================================================
\section{Results}
% ============================================================

% ------------------------------------------------------------
\subsection{A six-step pipeline identifies 728 high-confidence causal bQTL}
% ------------------------------------------------------------

We developed a filtering pipeline (Table~\ref{tab:pipeline}) that
progressively enriches for variants with unambiguous causal interpretation.
Starting from ${\sim}$237,000 MOA-seq binding peaks per hybrid, we first
identified 147,942 positions with allelic binding differences (bQTL). Of
these, 19,368 bQTL--gene pairs had a bQTL within 2\,kb of a gene's
transcription start site (TSS) and available ASE data.

The critical step is the genotype-stratified causal test. For each bQTL--gene
pair, we split the 19 genotyped NAM hybrids into two groups based on their
allele at the bQTL position, then compared the magnitude of allele-specific
expression (ASE) between groups using a Mann--Whitney $U$ test. Under the null
hypothesis that the bQTL has no effect on expression, the ASE distributions
should be identical regardless of genotype. We observed 1,403 significant
associations at FDR\,$<$\,0.05---a 4.85-fold enrichment over the ${\sim}$290
expected by chance ($p < 10^{-100}$, binomial test), confirming that the
genotype--ASE signal is overwhelmingly real.

To avoid LD confounding, we restricted to genes with exactly one significant
bQTL in their promoter (772 genes, 75\% of the significant set). Finally, we
required the bQTL to fall within at least one hybrid's MOA-seq peak file,
yielding the final set of 728 genes with high-confidence causal bQTL.

\begin{table}[htb]
\centering
\caption{\textbf{Filtering pipeline from bQTL to causal variants.}}
\label{tab:pipeline}
\begin{tabular}{clrr}
\toprule
\textbf{Step} & \textbf{Filter} & \textbf{N remaining} & \textbf{Drop} \\
\midrule
1 & All MOA binding peaks per hybrid & ${\sim}$237,000/hybrid & --- \\
2 & bQTL positions (allelic binding differences) & 147,942 & Non-variable \\
3 & bQTL within 2\,kb of gene TSS + ASE data & 19,368 pairs & Distal or no ASE \\
4 & Genotype $\to$ ASE causal test, FDR\,$<$\,0.05 & 1,403 pairs & Non-significant \\
5 & Single bQTL per gene (no LD ambiguity) & 772 genes & 257 multi-bQTL \\
6 & bQTL within MOA-seq peak & 728 genes & 44 not in peak \\
\bottomrule
\end{tabular}
\end{table}


% ------------------------------------------------------------
\subsection{Linkage disequilibrium and the single-bQTL design}
% ------------------------------------------------------------

Among the 1,403 significant bQTL--gene pairs, 75\% of genes have exactly one
significant bQTL, presenting no linkage concern. For the remaining 25\% with
multiple bQTL per gene, 62\% of within-gene variant pairs have identical
genotype vectors across the 19 founders (complete LD), making it impossible to
distinguish which variant is causal. Our single-bQTL filter removes this
ambiguity at the cost of excluding some real signals. This conservative design
ensures that every gene in the 728 set has an unambiguous causal variant.


% ------------------------------------------------------------
\subsection{Within binding peaks, disrupted motifs are non-redundant}
% ------------------------------------------------------------

A key concern with single-nucleotide regulatory variants is apparent
redundancy: if a promoter contains multiple copies of a TF motif, disrupting
one should have little effect. We addressed this by scanning the actual
MOA-seq binding peak sequence (extracted from the B73 reference genome) for
each of the 728 bQTL against all 259 PWMs from PlantTFDB~5.0 (36 TF
families, log-odds threshold\,=\,6.0). For each bQTL, we identified the
tightest peak containing the variant across all 25 hybrids, then counted how
many copies of each disrupted TF family's motif exist elsewhere within that
peak (Table~\ref{tab:redundancy}).

Of 728 bQTL, 600 (82.4\%) have at least one identifiable TF family motif at
the variant position. Among these, \textbf{47.8\%} (287/600) have at least
one disrupted family where the variant-overlapping motif is the \emph{sole
copy} of that family within the peak, and 13.0\% (78/600) have \emph{all}
disrupted families as sole copy. At the per-family-instance level, 28.2\%
(459/1,627) of all family-at-variant instances are sole copies.

Critically, sole-copy rates vary dramatically across TF families. Families
with longer, more information-rich motifs---AP2 (100\%), RAV (100\%), BES1
(86\%), Nin-like (85\%), CAMTA (70\%)---are almost always non-redundant
within the ${\sim}$160\,bp peak. In contrast, families with short,
degenerate motifs---C2H2 (10\%), ERF (10\%), MYB (9\%), Dof
(11\%)---maintain multiple copies even within narrow peaks. This family-level
variation suggests two distinct mechanisms: for high-specificity families,
the bQTL eliminates the only binding site; for low-specificity families, the
bQTL may reduce total binding affinity or shift the precise footprint
position without fully abolishing binding.

Importantly, sole-copy and redundant bQTL show comparable effect sizes
(median $|d|$\,=\,2.31 vs 2.44, Mann--Whitney $p$\,=\,0.15), indicating
that even redundant motif disruptions produce large ASE effects---consistent
with quantitative modulation of binding rather than binary on/off switching.

\begin{table}[htb]
\centering
\caption{\textbf{Motif redundancy within actual MOA-seq binding peaks.}
Sole-copy rates vary by TF family, with high-specificity families (BES1,
Nin-like) almost always non-redundant and low-specificity families (C2H2,
MYB) maintaining multiple copies even within ${\sim}$160\,bp peaks.}
\label{tab:redundancy}
\begin{tabular}{lrrr}
\toprule
\textbf{TF Family} & \textbf{At variant} & \textbf{Sole copy} & \textbf{\%} \\
\midrule
AP2       &   8 &   8 & 100 \\
RAV       &   6 &   6 & 100 \\
BES1      &   7 &   6 & 86 \\
Nin-like  &  27 &  23 & 85 \\
CAMTA     &  27 &  19 & 70 \\
HSF       &   9 &   6 & 67 \\
ARR-B     &  18 &  11 & 61 \\
bHLH      &  44 &  19 & 43 \\
LBD       &  51 &  21 & 41 \\
NAC       &  94 &  32 & 34 \\
G2-like   &  65 &  20 & 31 \\
MIKC-MADS &  92 &  24 & 26 \\
TALE      &  92 &  16 & 17 \\
ERF       & 123 &  12 & 10 \\
C2H2      & 136 &  14 & 10 \\
MYB       & 125 &  11 &  9 \\
\midrule
\textbf{All families} & \textbf{1,627} & \textbf{459} & \textbf{28.2} \\
\bottomrule
\end{tabular}
\end{table}


% ------------------------------------------------------------
\subsection{Variants do not just disrupt motifs---they also create new ones}
% ------------------------------------------------------------

A common assumption is that regulatory variants primarily \emph{disrupt}
existing TF binding sites. To test this, we performed a bidirectional motif
analysis: for each bQTL, we scanned both the B73 reference sequence and the
NAM variant sequence ($\pm$28\,bp around the variant) against all 259 PWMs
and classified each motif as disrupted (present in ref only), created
(present in alt only), or unchanged (Table~\ref{tab:creation},
Figure~\ref{fig:rewiring}).

Of 728 bQTL, 618 (84.9\%) alter at least one TF motif. Strikingly, only
24.2\% exclusively disrupt motifs; \textbf{21.2\% exclusively create new
motifs}, and \textbf{32.5\% simultaneously disrupt one family's motif while
creating another}. Effect sizes are indistinguishable
across all three categories (median $|d|$\,=\,2.38--2.44,
Figure~\ref{fig:rewiring_fx}A), confirming that motif creation produces
expression changes of equal magnitude to disruption.

At the family level, the net direction varies: NAC ($-$27), ERF ($-$20),
and TALE ($-$19) are strongly disruption-biased, while EIL ($+$11), WRKY
($+$6), and YABBY ($+$4) show net creation. MYB is perfectly balanced
(53 disrupted, 53 created). Zero self-switches were observed---when a
variant disrupts a family's motif, it never recreates a motif for the same
family, confirming that rewiring always involves a \emph{different} TF
(Figure~\ref{fig:rewiring}A). The most frequent transitions include C2H2
$\to$ TALE (6 bQTL), SBP $\to$ C2H2 (5), MYB-related $\to$ MYB (5), and
NAC $\to$ ERF (4).

\paragraph{Are rewired motifs sole copies within the peak?}
To determine whether rewiring represents a true regulatory switch rather
than a change among redundant copies, we repeated the analysis within actual
MOA-seq binding peaks using both the B73 and NAM allele sequences. Among
233 rewired bQTL, 22 (9.4\%) are \emph{clean switches}: all disrupted
families are sole copy in the reference peak, and all created families are
sole copy in the variant peak (Figure~\ref{fig:rewiring_fx}B). A further 30
(12.9\%) have at least one sole-copy family on each side. Clean switches
and partial switches show comparable effect sizes (median $|d|$\,=\,2.33
vs 2.42, $p$\,=\,0.57), suggesting that even partial rewiring is
functionally impactful.

\paragraph{Examples of clean sole-copy switches.}
Among the 22 clean switches, several involve genes at regulatory nodes:
\begin{itemize}
    \item \textbf{CAMTA $\to$ ERF} at phosphofructokinase
    (Zm00001eb048390, $|d|$\,=\,5.29): the sole CAMTA-family motif in the
    B73 peak is replaced by the sole ERF-family motif in the NAM peak,
    with the largest effect size among clean switches.
    \item \textbf{NAC $\to$ TALE/WRKY} at a proline-rich receptor kinase
    (Zm00001eb202030, $|d|$\,=\,2.62, drought DOWN): a sole NAC motif
    replaced by sole TALE and WRKY motifs at a drought-repressed signaling
    receptor.
    \item \textbf{HSF $\to$ HD-ZIP/WOX/YABBY} (Zm00001eb371180,
    $|d|$\,=\,2.42): the sole HSF motif replaced by three developmental
    TF family motifs in the variant peak.
\end{itemize}

Whether these switches redirect regulation between specific transcriptional
programs (e.g., stress versus development) or primarily alter local binding
affinity cannot be resolved from bulk tissue data (see Discussion). A
threshold sensitivity analysis (Methods) confirms that both the zero
self-switch finding and the equal effect sizes across categories are robust
across log-odds thresholds from 4.0 to 8.0.

\begin{table}[htb]
\centering
\caption{\textbf{Bidirectional motif changes at 728 causal bQTL.}
32.5\% of variants simultaneously disrupt one TF family's motif and create
another, rewiring rather than simply ablating TF occupancy.}
\label{tab:creation}
\begin{tabular}{lrrl}
\toprule
\textbf{Category} & \textbf{N} & \textbf{\%} & \textbf{Median $|d|$} \\
\midrule
Only disrupted   & 173 & 24.2 & 2.44 \\
Only created     & 152 & 21.2 & 2.42 \\
Both (rewired)   & 233 & 32.5 & 2.38 \\
Neither          & 158 & 22.1 & 2.36 \\
\midrule
\textbf{Total}   & \textbf{716} & \textbf{100} & 2.40 \\
\bottomrule
\end{tabular}
\end{table}

\begin{figure}[htb]
\centering
\includegraphics[width=\textwidth]{fig_motif_rewiring.pdf}
\caption{\textbf{TF family rewiring at causal bQTL.}
(A)~Transition heatmap: rows show disrupted TF families (B73 allele),
    columns show created families (NAM allele). No self-switches occur.
    C2H2~$\to$~TALE and SBP~$\to$~C2H2 are the most frequent transitions.
(B)~Category breakdown: rewired bQTL (32.5\%) are the largest category,
    with 22 clean sole-copy switches annotated.}
\label{fig:rewiring}
\end{figure}

\begin{figure}[htb]
\centering
\includegraphics[width=\textwidth]{fig_rewiring_effects.pdf}
\caption{\textbf{Effect sizes are independent of motif change category.}
(A)~All four categories show comparable effect sizes (median $|d| \approx 2.4$).
(B)~Among rewired bQTL, clean sole-copy switches ($n$\,=\,22) and partial
    switches ($n$\,=\,211) produce indistinguishable expression effects.}
\label{fig:rewiring_fx}
\end{figure}


% ------------------------------------------------------------
\subsection{Variant distribution across NAM hybrids}
% ------------------------------------------------------------

The 728 causal variants are distributed across 19 genotyped NAM hybrids with
striking uniformity (Figure~\ref{fig:hybrid_dist}). Each variant is carried by
a median of 7 hybrids (range: 3--16), with the distribution peaking at
intermediate frequencies (5--6 hybrids). No variant is fixed in all 19
hybrids, indicating that all causal variants are polymorphic in this
population. Conversely, each hybrid carries 264--308 causal variants
(36--42\% of the 728), a remarkably narrow range suggesting that the
\emph{total regulatory variant burden} is approximately equal across diverse
maize germplasm.

Rare variants (carried by $\leq$4 hybrids) show slightly larger ASE effect
sizes than common variants ($\rho = -0.11$, $p = 0.003$), consistent with
purifying selection acting more strongly against large-effect regulatory
mutations. The effect size (Cohen's $d$) for the causal test ranges from
1.0 to 14.4 (median 2.4), reflecting substantial allelic expression
differences.

\begin{figure}[htb]
\centering
\includegraphics[width=\textwidth]{fig_causal_bqtl_hybrid_distribution.pdf}
\caption{\textbf{Distribution of 728 causal bQTL across NAM hybrids.}
(A)~Histogram of the number of hybrids carrying each variant.
(B)~ASE effect size (|Cohen's $d$|) versus variant frequency, with
    weak negative correlation (Spearman $\rho = -0.11$, $p = 0.003$).
(C)~Number of causal variants per hybrid (264--308 per hybrid).
(D)~Cumulative distribution of variants by number of carriers.}
\label{fig:hybrid_dist}
\end{figure}


% ------------------------------------------------------------
\subsection{Target genes encode the regulatory layer, not metabolic enzymes}
% ------------------------------------------------------------

Functional characterization of the 728 target genes reveals a clear bias
toward regulatory and signaling functions (Figure~\ref{fig:signaling}A):
kinases/phosphorylation (76 genes, 10.4\%), chaperones/HSPs (63, 8.7\%),
biosynthetic enzymes (39, 5.4\%), transcription factors (32, 4.4\%),
translation machinery (29, 4.0\%), signal transduction (26, 3.6\%),
and transporters (17, 2.3\%).

Pathway enrichment analysis (Fisher's exact test against genome background)
identified five significantly enriched KEGG pathways: pentose phosphate
(OR\,=\,3.37, $p$\,=\,0.012), steroid biosynthesis (OR\,=\,3.60,
$p$\,=\,0.016), amino sugar/nucleotide sugar metabolism (OR\,=\,2.11,
$p$\,=\,0.021), ubiquitin-mediated proteolysis (OR\,=\,2.14,
$p$\,=\,0.015), and pantothenate/CoA biosynthesis (OR\,=\,4.32,
$p$\,=\,0.018). Enriched GO terms include pentose metabolic process
(OR\,=\,19.2, $p$\,=\,$1.6 \times 10^{-4}$) and regulation of hormone
biosynthesis (OR\,=\,5.1, $p$\,=\,0.011).

Notably, large metabolic gene clusters (benzoxazinoid, terpene synthase) are
largely absent from the 728 set---only Bx9 (UDP-glucosyltransferase) and a
single flavonoid 3'-monooxygenase are present. This confirms that natural
cis-regulatory variation predominantly targets the \emph{regulatory layer}:
the variant does not change an enzyme directly, but modulates how much a
regulatory gene is expressed from one allele, which cascades to downstream
effects.

We classified all 728 proteins into cellular functional categories
(Figure~\ref{fig:signaling}): signal transduction ($n$\,=\,129) is the
largest, followed by protein quality control ($n$\,=\,64), metabolism
($n$\,=\,50), translation ($n$\,=\,31), and transcription factors
($n$\,=\,28). Receptors/sensors ($n$\,=\,21), membrane/transport
($n$\,=\,24), and hormone signaling ($n$\,=\,7) complete the picture. Effect
sizes are comparable across all layers (median $|d|$\,=\,2.4--2.9),
indicating no preferential impact on any particular level of the signaling
hierarchy.


\begin{figure}[htb]
\centering
\includegraphics[width=\textwidth]{fig_signaling_layers.pdf}
\caption{\textbf{Functional architecture of 728 causal bQTL target genes.}
(A)~Distribution across functional categories, with mean effect size annotated.
(B)~TF motif $\times$ functional category heatmap showing that all TF families
    regulate targets across all layers.
(C)~Effect size distributions are comparable across functional categories.
(D)~Drought response (|log2FC|\,$>$\,1) by functional category; signal transduction
    and protein quality control show the strongest response.}
\label{fig:signaling}
\end{figure}


% ------------------------------------------------------------
\subsection{Gene $\times$ TF family heatmap reveals co-regulation structure}
% ------------------------------------------------------------

To map which TF family motifs are disrupted at each gene's bQTL, we scanned
all 728 variant positions against 259 PWMs. Of 728 genes, 662 (90.9\%) have
at least one identifiable TF family motif at the variant position, covering 35
families. The most frequently disrupted families are C2H2 and MYB (each 160
genes), ERF (151), NAC (118), TALE (115), and MIKC-MADS (107).

Hierarchical clustering of the gene $\times$ TF family binary matrix
(Figure~\ref{fig:heatmap}) reveals that most genes have 2--5 TF families
disrupted at their bQTL position, reflecting the overlapping nature of TF
binding motifs. Only 134 genes (18.4\%) have a single exclusive TF family
disrupted. Among these, ERF and MYB each exclusively control 20 target genes,
followed by NAC (9), Dof (8), and SBP (7).

The Jaccard overlap between TF family target sets
(Figure~\ref{fig:coreg}B) shows moderate overlap ($J$\,=\,0.2--0.4),
confirming that the same genes often have motifs for multiple families
disrupted at a single variant position. This is consistent with the dense
packing of TF motifs in regulatory regions and the degenerate nature of
PWM matching.

\begin{figure}[htb]
\centering
\includegraphics[width=0.95\textwidth]{fig_gene_motif_heatmap.pdf}
\caption{\textbf{Gene $\times$ TF family heatmap.}
662 genes with identifiable disrupted motifs (rows) $\times$ 35 TF families
(columns). Blue indicates the TF family's motif is disrupted at the bQTL
position. Genes are hierarchically clustered by Jaccard distance; bar chart
below shows total genes per family.}
\label{fig:heatmap}
\end{figure}


% ------------------------------------------------------------
\subsection{Co-regulation patterns and TF family convergence}
% ------------------------------------------------------------

Despite the diversity of disrupted TF families, all families regulate
strikingly similar target gene profiles (Figure~\ref{fig:signaling}B). The
top target functions for every major TF family are dominated by signal
transduction, protein quality control, and transcription factors. No family
is enriched for a specific functional category (Figure~\ref{fig:signaling}B).

The coefficient of variation (CV) in mean effect size across TF families is
only 5.5\%, confirming that TF identity has minimal influence on the
magnitude of allelic expression differences. This convergence supports the
model that \emph{chromatin accessibility}---whether the variant is in open
chromatin and physically accessible to any TF---is the primary determinant of
functional impact, not which specific TF family's motif is disrupted
\citep{engelhorn2025}.

\begin{figure}[htb]
\centering
\includegraphics[width=\textwidth]{fig_coregulation_network.pdf}
\caption{\textbf{Co-regulation patterns among 728 causal bQTL.}
(A)~Genes exclusively regulated by a single TF family (sole-motif groups).
    ERF and MYB each control 20 genes exclusively.
(B)~Jaccard similarity between TF family target gene sets. Moderate overlap
    ($J$\,=\,0.2--0.4) reflects shared motif occurrence at variant sites.}
\label{fig:coreg}
\end{figure}


% ------------------------------------------------------------
\subsection{TF cascade network: 34 transcription factor targets
amplify regulatory variation}
% ------------------------------------------------------------

Among the 728 target genes, 34 are themselves transcription factors, including
MYB30, MYB44, WRKY50, WRKY51, NAC67, NAC82, bHLH68, bHLH112, and bHLH128.
These create regulatory cascades where a single upstream motif disruption can
propagate to many downstream genes (Figure~\ref{fig:cascade}).

The most powerful cascade nodes are:
\begin{itemize}
    \item \textbf{MYB30} ($|d|$\,=\,5.31, drought DOWN): disrupted by a MYB
    motif variant, MYB30 regulates 96 downstream genes in the 728 set. Its
    strong drought repression suggests it normally promotes growth-related
    targets that are silenced under stress.
    \item \textbf{MYB44} ($|d|$\,=\,2.28, drought UP): disrupted by a GATA
    motif variant, regulates 96 downstream genes. Its drought induction and
    known role in ABA signaling make it a key stress-responsive cascade node.
    \item \textbf{WRKY51} ($|d|$\,=\,2.03, drought UP): disrupted by
    GATA/HD-ZIP/WOX/YABBY motifs, a known stress-responsive TF.
    \item \textbf{bHLH112} ($|d|$\,=\,2.51, drought UP): disrupted by a TALE
    motif, involved in drought and salt tolerance.
\end{itemize}

These cascade TFs have opposing drought responses (MYB30 DOWN vs MYB44 UP),
suggesting that natural variants in their promoters can rewire the
drought response by simultaneously reducing one branch of the regulatory
network while activating another.

\begin{figure}[htb]
\centering
\includegraphics[width=\textwidth]{fig_tf_cascade_network.pdf}
\caption{\textbf{TF cascade network among 728 causal bQTL targets.}
34 TF genes (middle layer) are themselves targets of bQTL that disrupt
upstream motifs (top layer). Each TF in turn regulates downstream target genes
(bottom layer), creating regulatory amplification cascades.}
\label{fig:cascade}
\end{figure}


% ------------------------------------------------------------
\subsection{Drought integration: 27\% of target genes are
drought-responsive}
% ------------------------------------------------------------

Cross-referencing with drought expression data from \citet{engelhorn2025}
reveals that 16.5\% of target genes are drought-upregulated
(log2FC\,$>$\,1) and 10.7\% are drought-downregulated (log2FC\,$<$\,$-1$),
with a slight overall induction bias (mean log2FC\,=\,$+$0.206). Among
functional categories, signal transduction and protein quality control show the
strongest drought response asymmetry (Figure~\ref{fig:signaling}D).

Notable drought-responsive targets include:
\begin{itemize}
    \item \textbf{CIPK31} (CBL-interacting protein kinase): log2FC\,=\,$+$1.38,
    $|d|$\,=\,4.11---a calcium-signaling kinase with one of the largest effect
    sizes among drought-responsive genes.
    \item \textbf{PYR1} (ABA receptor): $|d|$\,=\,2.44---the bQTL variant
    affects expression of the core ABA perception machinery.
    \item \textbf{ARF7} (auxin response factor): $|d|$\,=\,3.47---linking
    cis-regulatory variation to auxin-mediated growth responses.
    \item \textbf{Phosphoenolpyruvate carboxylase 1}: log2FC\,=\,$+$3.58,
    $|d|$\,=\,2.48---the most drought-induced metabolic gene, connecting
    carbon metabolism to stress adaptation.
\end{itemize}

These examples illustrate how causal bQTL can be traced from a single
nucleotide change through TF binding disruption to altered expression of
defined stress signaling components.


% ------------------------------------------------------------
\subsection{Known signaling modules within the 728 set}
% ------------------------------------------------------------

Several characterized plant signaling pathways are represented by multiple
members in the 728 set:

\paragraph{Auxin signaling.} Five genes including IAA5 (auxin-responsive
protein), ARF7 (auxin response factor), and auxin-induced protein 15A are
regulated by variants disrupting C2H2, TALE, Dof, and TCP motifs.

\paragraph{ABA signaling.} Fourteen genes including the PYR1 ABA receptor,
regulated by variants in ERF, TCP, and CPP motifs. This module connects
genotype to the core drought perception pathway.

\paragraph{HSP/chaperone network.} Forty-eight genes form the largest coherent
module, with disrupted motifs spanning many TF families. This convergence
suggests that chaperone expression is broadly sensitive to regulatory
variation, consistent with chaperones acting as phenotypic buffers
\citep{rutherford1998}.

\paragraph{Ubiquitin--proteasome system.} Thirty-two genes involved in protein
turnover, enriched in the 728 set (OR\,=\,2.14, $p$\,=\,0.015). Targeted
protein degradation is a central mechanism for rapid signal transduction in
plants.


% ============================================================
\section{Discussion}
% ============================================================

Our analysis of 728 high-confidence causal bQTL provides several insights
into the architecture of natural gene expression variation in maize.

\paragraph{The regulatory layer is the primary target.}
The 728 causal variants overwhelmingly target signaling and regulatory
proteins (kinases, TFs, chaperones, transporters) rather than metabolic
enzymes or structural proteins. This is consistent with the omnigenic model
\citep{boyle2017}, where core gene expression is modulated by a large network
of peripheral regulatory genes. Natural selection appears to act on the
\emph{wiring} of the regulatory network rather than on the enzymes themselves.

\paragraph{Chromatin accessibility gates functional variation.}
The finding that sole-copy rates range from 9\% (MYB) to 100\% (AP2)
within the ${\sim}$160\,bp binding peak---yet both sole-copy and redundant
bQTL produce comparable effect sizes ($|d|$\,=\,2.3 vs 2.4)---supports a
model of quantitative modulation rather than binary on/off switching.
Combined with the observation that all TF families regulate similar target
profiles, this suggests that the primary determinant of whether a variant
``matters'' is not which TF it affects, but whether the variant is in open,
accessible chromatin. This echoes recent findings in human genetics where
chromatin accessibility is the strongest predictor of GWAS variant
functionality \citep{finucane2015}.

\paragraph{Variants rewire TF occupancy, not just disrupt it.}
The finding that 32.5\% of bQTL simultaneously disrupt one TF family's motif
and create another challenges the standard ``loss-of-binding'' model of
regulatory variation. However, three observations suggest that the TF family
\emph{label} of the switch may be less important than the change in local
binding energy: (1)~effect sizes are identical across all categories
(disrupted, created, rewired, neither); (2)~switching to a more
abundantly expressed TF family does not predict upregulation (Spearman
$\rho$\,=\,0.06, $p$\,=\,0.36); and (3)~growth/developmental versus
stress/defense switches are symmetric (44\% vs 56\%, binomial
$p$\,=\,0.45). These results are consistent with a model where the variant
alters the \emph{physical properties} of the binding site (DNA shape,
flexibility, binding free energy) in a way that affects occupancy by
whichever TF happens to be present, rather than redirecting regulation from
one specific transcriptional program to another.

A critical caveat is that our data derive from bulk leaf tissue, averaging
signal across mesophyll, epidermis, bundle sheath, guard cells, and
vasculature. Different cell types express different TF repertoires: a CAMTA
motif may be functionally decisive in guard cells (where calcium signaling
dominates), while an ERF motif may matter in mesophyll (ethylene response).
In bulk, each TF family is the ``important'' one in some subset of cells,
producing the appearance that all families are equivalent. Cell-type-resolved
MOA-seq or single-cell chromatin accessibility data would likely reveal
cell-type-specific TF hierarchies where family identity does predict effect
direction---a prediction that is now directly testable.

\paragraph{Regulatory cascades amplify single-variant effects.}
The 34 TF genes that are themselves bQTL targets create an amplification
layer: a variant in MYB30's promoter does not just change MYB30 expression,
but potentially affects 96 downstream genes. Importantly, cascade TFs with
opposing drought responses (MYB30 DOWN, MYB44 UP) suggest that natural
variation can simultaneously activate and repress different branches of the
stress response network.

\paragraph{Implications for precision breeding.}
Each of the 728 variants represents a defined molecular switch: introducing
or restoring a single nucleotide quantitatively modulates TF binding within
the peak, causally shifting expression with a quantified effect size (Cohen's
$d$). The median variant is shared by 7 hybrids, making effects testable
across genetic backgrounds. Combined with the interactive gene atlas
(Supplementary File~S1), this provides a prioritized catalog for CRISPR-based
promoter engineering, where the predicted effect direction and magnitude are
known \emph{a priori}.

\paragraph{Limitations.}
Our analysis uses 19 genotyped hybrids, limiting statistical power for
rare variants. The single-bQTL design excludes 25\% of significant genes
with multiple bQTL, potentially missing combinatorial regulatory effects.
TF family assignments depend on PWM quality, and motifs for some families
may be underrepresented in the PlantTFDB~5.0 database. The phenotypic
consequences of these regulatory variants cannot be resolved at the
individual-gene level with $N$\,=\,19 hybrids and will require larger
populations or direct functional validation.


% ============================================================
\section{Methods}
% ============================================================

\subsection{Data sources}
All bQTL data derive from \citet{engelhorn2025}, comprising MOA-seq binding
peaks and allele-specific expression measurements from 25~maize F1 hybrids
(B73 $\times$ NAM founders) under well-watered conditions. Genotype data
from assembled founder genomes are available for 19 of 25 hybrids (excluding
A188, A619, CML103, Mo17, W22). Gene annotations use the B73 NAM5 reference
genome (Zm-B73-REFERENCE-NAM-5.0), with functional annotations from EntAP
and EggNOG~5.0.

\subsection{Genotype-stratified causal test}
For each bQTL--gene pair, we extracted the genotype (reference or variant
allele) at the bQTL position across 19 hybrids. We then compared the
distribution of ASE values (log2(B73/NAM) expression ratios) between the
variant-carrying and reference groups using a two-sided Mann--Whitney $U$
test. Effect sizes were quantified as Cohen's $d$. $P$-values were corrected
for multiple testing using the Benjamini--Hochberg procedure.

\subsection{Motif scanning and bidirectional analysis}
TF binding motifs were scanned using 259 PWMs from PlantTFDB~5.0 (36 TF
families) with a log-odds threshold of 6.0. The variant-overlap constraint
required motif hits to physically overlap the variant nucleotide position
($\pm$28\,bp, the maximum PWM width). Motif redundancy was assessed within
actual MOA-seq binding peaks (tightest peak per bQTL across all 25 hybrids).
For bidirectional analysis, both the B73 reference and NAM variant sequences
were scanned; motifs present only in the reference were classified as
disrupted, those present only in the variant as created.

\subsection{Threshold sensitivity analysis}
To assess robustness, the bidirectional motif analysis was repeated at seven
log-odds thresholds (4.0, 5.0, 5.5, 6.0, 6.5, 7.0, 8.0). As expected, the
rewired fraction varies with threshold (18.5\%--59.9\%), reflecting more
permissive motif calling at lower thresholds. However, three key findings
are stable across all thresholds: (1)~zero self-switches at every threshold;
(2)~motif creation remains substantial (14.6\%--21.3\% of bQTL exclusively
create motifs); and (3)~effect sizes are indistinguishable across categories
at every threshold (median $|d|$ range 2.29--2.54).

\subsection{Functional category classification}
All 728 genes were classified into cellular functional categories based on
keyword matching against EntAP protein descriptions, EggNOG COG
descriptions, UniProt GO terms, and KEGG pathway annotations. Layers include
receptor/sensor, membrane/transport, signal transduction, protein quality
control, transcription factor, chromatin/epigenetic, translation, metabolism,
hormone signaling, and redox/stress response.

\subsection{Interactive atlas}
The Maize Causal bQTL Atlas (\url{https://mbanf.github.io/causal-bqtl-atlas/})
is an interactive HTML visualization built with D3.js~v7. It includes:
a force-directed layered cell diagram of all 728 genes colored by
regulatory layer; rich tooltips with gene description, effect size,
drought response, disrupted and created TF motifs, motif rewiring
category, sole-copy status within MOA-seq peaks, and hybrid carriers;
filterable by motif rewiring category (rewired, only disrupted, only
created, neither), regulatory layer, disrupted and created TF family
(union filter), drought category, and specific hybrid; and an optional
OpenAI API integration for AI-powered biological summaries of selected
gene sets.


% ============================================================
\section{Data Availability}
% ============================================================

All bQTL data are from \citet{engelhorn2025}: SRA accession PRJNA1101486,
GEO accession GSE294039. The 728-gene atlas (gene\_atlas\_728.csv),
interactive visualization (gene\_atlas\_728\_interactive.html), and analysis
scripts are available at \url{https://github.com/mbanf/causal-bqtl-atlas}.


% ============================================================
% REFERENCES
% ============================================================

\bibliographystyle{unsrtnat}
\begin{thebibliography}{99}

\bibitem[Engelhorn et al.(2025)]{engelhorn2025}
Engelhorn, J., Snodgrass, S.J., Banf, M., et al. (2025).
A maize leaf pan-cistrome catalogs the cis-regulatory variation underlying
functional trait heritability.
\emph{Nature Genetics}, s41588-025-02246-7.

\bibitem[Boyle et al.(2017)]{boyle2017}
Boyle, E.A., Li, Y.I., Pritchard, J.K. (2017).
An expanded view of complex traits: from polygenic to omnigenic.
\emph{Cell}, 169(7), 1177--1186.

\bibitem[Finucane et al.(2015)]{finucane2015}
Finucane, H.K., et al. (2015).
Partitioning heritability by functional annotation using genome-wide
association summary statistics.
\emph{Nature Genetics}, 47(11), 1228--1235.

\bibitem[Rutherford and Lindquist(1998)]{rutherford1998}
Rutherford, S.L. and Lindquist, S. (1998).
Hsp90 as a capacitor for morphological evolution.
\emph{Nature}, 396(6709), 336--342.

\bibitem[Segal et al.(2008)]{segal2008}
Segal, E. and Widom, J. (2008).
From DNA sequence to transcriptional behaviour: a quantitative approach.
\emph{Nature Reviews Genetics}, 10(7), 443--456.

\bibitem[Levo and Segal(2012)]{levo2012}
Levo, M. and Segal, E. (2012).
In pursuit of design principles of regulatory sequences.
\emph{Nature Reviews Genetics}, 13(11), 786--798.

\end{thebibliography}


\end{document}
